% !TeX program = xelatex
% !TeX encoding = UTF-8
% Software Architecture Document -- adapted from the lecturer's
% simplified RUP Software Architecture Document Word template.
%
% Folder layout (keep main.tex and metadata.tex together):
%   main.tex
%   metadata.tex
%   sections/
%     01-Introduction.tex
%     02-Architectural Goals and Constraints.tex
%     03-Use-Case Model.tex
%     04-Logical View.tex
%     05-Deployment.tex
%     06-Implementation View.tex
%
% Section files own their headings and fillable content.
% Build from this directory, repeating until the TOC/references stabilize:
%   xelatex -interaction=nonstopmode -halt-on-error -output-directory=build main.tex
% Create build/ first; keep all generated files there.

\documentclass[10pt,letterpaper]{article}
% Project and document metadata only; layout belongs in main.tex.
% Project/team/organization/author names are retained from the PA2 metadata.
% Team details are carried over from the PA2 project documentation.
% Escape LaTeX special characters in values: \&, \%, \_, \#.

\newcommand{\ProjectName}{MEDICATION SAFETY ASSISTANT}
\newcommand{\GroupName}{Medication Safety Assistant Team}
\newcommand{\GroupNumber}{Group 02}
\newcommand{\OrganizationName}{University of Science, VNU-HCM}
\newcommand{\DocumentAuthors}{Cao Hoàng Long, Hoàng Yến Trang, Lỗ Chí Thành}

\newcommand{\DocumentTitle}{Software Architecture Document}
\newcommand{\DocumentPhase}{PA3}
\newcommand{\DocumentIdentifier}{PA3-Group02-SAD}
\newcommand{\DocumentVersion}{1.0}
\newcommand{\DocumentDate}{11 November 2026}
\newcommand{\DocumentYear}{2026}
\newcommand{\ConfidentialityText}{Confidential}

\newcommand{\RevisionDate}{11 Nov 2026}
\newcommand{\RevisionVersion}{\DocumentVersion}
\newcommand{\RevisionDescription}{Initial PA3 Software Architecture Document template based on the instructor's simplified RUP format}
\newcommand{\RevisionAuthor}{Group 02}

% -------------------- Fonts and layout -------------------------
\usepackage{fontspec}
\defaultfontfeatures{Ligatures=TeX}
\IfFontExistsTF{Times New Roman}{%
  \setmainfont{Times New Roman}
}{%
  \IfFontExistsTF{Tinos}{%
    \setmainfont{Tinos}
  }{%
    \setmainfont{Liberation Serif}
  }%
}
\IfFontExistsTF{Arial}{\setsansfont{Arial}}{%
  \IfFontExistsTF{TeX Gyre Heros}{\setsansfont{TeX Gyre Heros}}{\setsansfont{Liberation Sans}}
}
% The instructor's RUP document uses US Letter; typography follows the PA2 Vision style.
\usepackage[letterpaper,margin=1in,headheight=34pt,headsep=16pt,
            footskip=28pt]{geometry}
\usepackage{setspace}
\setstretch{1.08}
\usepackage{microtype}
\usepackage{array,tabularx,longtable,booktabs}
\usepackage{graphicx}
\graphicspath{{images/}{figures/}}
\usepackage{float}
\usepackage{caption}
\usepackage{pdflscape}
\usepackage{enumitem}
\usepackage{titlesec}
\usepackage{fancyhdr}
\usepackage{lastpage}
\usepackage{xcolor}
\usepackage{needspace}
\usepackage{tocloft}
\usepackage[unicode,hidelinks]{hyperref}

\setlength{\parindent}{0pt}
\setlength{\parskip}{5pt}
\setlength{\tabcolsep}{5pt}
\renewcommand{\arraystretch}{1.23}
\emergencystretch=2em
\setcounter{secnumdepth}{3}
\setcounter{tocdepth}{3}
\setlist[itemize]{topsep=3pt,itemsep=2pt,parsep=0pt}
\renewcommand{\contentsname}{Table of Contents}
\titleformat{\section}{\normalfont\bfseries\fontsize{12}{15}\selectfont}{\thesection.}{0.6em}{}
\titleformat{\subsection}{\normalfont\bfseries\normalsize}{\thesubsection}{0.6em}{}
\titleformat{\subsubsection}{\normalfont\itshape\normalsize}{\thesubsubsection}{0.6em}{}
\titlespacing*{\section}{0pt}{6pt}{3pt}
\titlespacing*{\subsection}{0pt}{6pt}{3pt}
\renewcommand{\cfttoctitlefont}{\hfill\normalfont\bfseries\fontsize{17}{21}\selectfont}
\renewcommand{\cftaftertoctitle}{\hfill\mbox{}}
\setlength{\cftbeforetoctitleskip}{0pt}
\setlength{\cftaftertoctitleskip}{16pt}
\renewcommand{\cftsecfont}{\normalfont}
\renewcommand{\cftsecpagefont}{\normalfont}
\renewcommand{\cftsecaftersnum}{.}
\renewcommand{\cftsecleader}{\cftdotfill{\cftdotsep}}
\renewcommand{\cftsubsecleader}{\cftdotfill{\cftdotsep}}
\setlength{\cftbeforesecskip}{7pt}
\setlength{\cftsecnumwidth}{1.7em}
\setlength{\cftsubsecindent}{1.7em}
\setlength{\cftsubsecnumwidth}{3em}
\tocloftpagestyle{fancy}

% Compatible with section files containing figures, tables, diagrams.
\newcolumntype{L}[1]{>{\raggedright\arraybackslash}p{#1}}
\newcolumntype{Y}{>{\raggedright\arraybackslash}X}
\newcommand{\sourcefile}[1]{\nolinkurl{#1}}
\newcommand{\db}[1]{\nolinkurl{#1}}
\newcommand{\status}[1]{\textsc{#1}}
% Blue italic instructions follow the instructor's InfoBlue guidance style.
% Replace each TODO with authored content and remove the instruction wrapper.
\DeclareRobustCommand{\TODO}[1]{{\color{blue}\itshape TODO: #1}}
\newcommand{\diagramplaceholder}[4][3cm]{%
  \begin{figure}[H]
    \centering
    \fbox{\begin{minipage}[c][#1][c]{\dimexpr\linewidth-2\fboxsep-2\fboxrule\relax}
      \centering\TODO{#2}
    \end{minipage}}
    \caption{\TODO{#3}}\label{#4}
  \end{figure}%
}
\newcommand{\frontmattertitle}[1]{%
  \begin{center}{\normalfont\bfseries\fontsize{16}{20}\selectfont #1\par}\end{center}%
}
\newcommand{\portraitdiagram}[3]{%
  \begin{figure}[p]
    \centering
    \includegraphics[width=\textwidth,height=.78\textheight,
                     keepaspectratio]{#1}
    \caption{#2}\label{#3}
  \end{figure}
  \clearpage
}
\newcommand{\landscapediagram}[3]{%
  \clearpage
  \begin{landscape}
    \begin{figure}[p]
      \centering
      \includegraphics[width=\linewidth,height=.82\textheight,
                       keepaspectratio]{#1}
      \caption{#2}\label{#3}
    \end{figure}
  \end{landscape}
  \clearpage
}

% --------- Clean Vision-style running header/footer ----------
% Do not use a tabular/table here: this is a running header, not a data table.
\pagestyle{fancy}
\fancyhf{}
\fancyhead[L]{%
  \begin{minipage}[b]{.73\textwidth}
    \raggedright\fontsize{9}{11}\selectfont
    \ProjectName\\[-1pt]
    {\fontsize{8}{10}\selectfont\DocumentTitle\enspace(\DocumentIdentifier)}
  \end{minipage}%
}
\fancyhead[R]{%
  \begin{minipage}[b]{.24\textwidth}
    \raggedleft\fontsize{9}{11}\selectfont
    Version \DocumentVersion\\[-1pt]
    {\fontsize{8}{10}\selectfont\DocumentDate}
  \end{minipage}%
}
\fancyfoot[L]{\fontsize{8.5}{10}\selectfont\OrganizationName{} -- Course Project}
\fancyfoot[C]{\fontsize{8}{10}\selectfont\ConfidentialityText}
\fancyfoot[R]{\fontsize{8.5}{10}\selectfont Page \thepage\ of \pageref{LastPage}}
\renewcommand{\headrulewidth}{0.45pt}
\renewcommand{\footrulewidth}{0.45pt}
% Like the PA2 Vision cover, the title page has no running header/footer.
\fancypagestyle{sadcover}{\fancyhf{}\renewcommand{\headrulewidth}{0pt}%
  \renewcommand{\footrulewidth}{0pt}}

\hypersetup{
  pdftitle={\DocumentTitle\ - \ProjectName},
  pdfauthor={\DocumentAuthors},
  pdfsubject={\DocumentPhase\ Software Architecture}
}

\begin{document}

% ------------------ PAGE 1: Vision-style cover ------------------
\thispagestyle{sadcover}
\begin{center}
  \vspace*{.35in}
  {\bfseries\fontsize{13}{17}\selectfont \ProjectName\par}
  \vspace*{1.50in}
  {\bfseries\fontsize{22}{27}\selectfont \DocumentTitle\par}
  \vspace{8pt}
  {\fontsize{12}{16}\selectfont (Small Project)\par}
  \vspace{.92in}
  {\fontsize{12}{15}\selectfont Version \DocumentVersion\par}
  \vfill
  {\fontsize{10}{14}\selectfont
    \begin{tabular}{@{}r@{\quad}l@{}}
      \textbf{Document identifier:} & \DocumentIdentifier\\
      \textbf{Prepared by:} & \GroupNumber{} -- \GroupName\\
      \textbf{Authors:} & \DocumentAuthors\\
      \textbf{Organization:} & \OrganizationName\\
      \textbf{Date:} & \DocumentDate
    \end{tabular}\par}
  \vspace*{1.05in}
\end{center}
\clearpage

% ---------------- PAGE 2: Revision History ---------------------
\frontmattertitle{Revision History}
\vspace{7pt}
{\renewcommand{\arraystretch}{1.28}%
\begin{tabularx}{\textwidth}{|L{2.10cm}|L{1.30cm}|Y|L{2.60cm}|}
  \hline
  \textbf{Date} & \textbf{Version} & \textbf{Description} & \textbf{Author}\\
  \hline
  \RevisionDate & \RevisionVersion & \RevisionDescription & \RevisionAuthor\\
  \hline
\end{tabularx}}
\clearpage


% ---------------- Table of Contents --------------------
\tableofcontents
\clearpage

% ---------------- List of Figures --------------------
\phantomsection
\addcontentsline{toc}{section}{List of Figures}
\listoffigures
\clearpage

% ---------------- List of Tables --------------------
\phantomsection
\addcontentsline{toc}{section}{List of Tables}
\listoftables
\clearpage

% ---------------- MAIN BODY: lecturer's exact sections ---------
\frontmattertitle{\DocumentTitle}

% 1. Overview, purpose, scope, terminology, references.

\section{Introduction}
\label{sec:introduction}

\noindent\textbf{Purpose.}
This document describes the software architecture of the
\emph{Medication Safety Assistant}, including its main components,
responsibilities, interfaces, and relationships. It serves as
a technical reference for the project team and course reviewers.

\medskip
\noindent\textbf{Scope.}
The Medication Safety Assistant is a Web-based research prototype
supporting medication intake, identification, safety assessment,
simulated healthcare professional review, and medication reminders.

This document covers architectural goals and constraints,
the use-case model, and the logical view. The system provides
decision support only and does not replace professional medical
judgment. Deployment and implementation views are deferred
as permitted by the PA3 requirements.

\medskip
\noindent\textbf{Definitions, Acronyms and Abbreviations.}

\noindent
\begin{tabularx}{\textwidth}{|L{2.5cm}|Y|}
    \hline
    \textbf{Term} & \textbf{Definition} \\ \hline
    HCP & Healthcare Professional \\ \hline
    DDI & Drug--Drug Interaction \\ \hline
    OCR & Optical Character Recognition \\ \hline
    NLP & Natural Language Processing \\ \hline
    RBAC & Role-Based Access Control \\ \hline
\end{tabularx}

\medskip
\noindent\textbf{References.}

\noindent
\begin{tabularx}{\textwidth}{|L{1.5cm}|Y|L{2.8cm}|}
    \hline
    \textbf{ID} & \textbf{Document} & \textbf{Version / Date} \\ \hline
    R1 & Vision Document (PA2-Group02-VisionDocument)
       & 1.1 / 01 Nov 2026 \\ \hline
    R2 & Use-Case Specification (PA2-Group02-UseCaseSpecification)
       & 1.0 / 01 Nov 2026 \\ \hline
    R3 & Software Development Plan (PA2-Group02-ProjectPlan)
       & 1.1 / 01 Nov 2026 \\ \hline
\end{tabularx}

\medskip
\noindent\textbf{Overview.}
Section 2 presents architectural goals and constraints.
Section 3 presents the use-case model.
Section 4 describes the logical architecture and its components.
Sections 5 and 6 are reserved for future deployment and
implementation details.


% 2. Key non-functional requirements and constraints from Vision.

\section{Architectural Goals and Constraints}
\label{sec:architectural-goals}

\medskip
\noindent\textbf{Architectural Goals}

\begin{table}[!htbp]
\centering
\caption{Architectural Goals and Their Impacts}
\label{tab:architectural-goals}

\begin{tabularx}{\textwidth}{|L{3cm}|L{2cm}|Y|}
\hline
\textbf{Goal} & \textbf{Source} &
\textbf{Architectural Impact} \\ \hline

Safety and Reliability & NFR-02, NFR-03 &
Confirm medication input, handle incomplete assessments,
preserve severe warnings, and support fallback input.
\\ \hline

Security and Privacy & NFR-04 &
Enforce RBAC and patient-authorized case sharing.
\\ \hline

Data Integrity and Traceability & NFR-07, NFR-10 &
Preserve medication history, assessment results,
review decisions, and source references.
\\ \hline

Performance and Usability & NFR-01, NFR-05 &
Complete at least 90\% of standard safety checks within
10 seconds under normal conditions; provide clear
warnings and status feedback.
\\ \hline

Maintainability and Extensibility & NFR-08 &
Separate the five main modules with clearly defined
responsibilities.
\\ \hline

\end{tabularx}
\end{table}

\medskip
\noindent\textbf{Architectural Constraints}

\begin{table}[!htbp]
\centering
\caption{Architectural Constraints and Their Impacts}
\label{tab:architectural-constraints}

\begin{tabularx}{\textwidth}{|L{3cm}|L{2cm}|Y|}
\hline
\textbf{Constraint} & \textbf{Source} &
\textbf{Architectural Impact} \\ \hline

Web Platform & NFR-06 &
Support major desktop and mobile browsers;
provide manual input and in-app schedule viewing
when browser capabilities are unavailable.
\\ \hline

Medication Safety Rules & NFR-03, PF-07 &
Activate only verified, confirmed schedules.
Hold new reminders for medicines with severe
interactions pending HCP review; never
automatically stop existing reminders or treatment.
\\ \hline

Research Scope and Data Sources & NFR-04, NFR-10 &
Use synthetic or de-identified data and simulated HCP
review. Respect data-source permissions and
mark unsupported checks as unassessed.
\\ \hline

\end{tabularx}
\end{table}


% 3. Use-case diagram(s) from the Use-Case Specification document.

\section{Use-Case Model}
\label{sec:use-case-model}

This section presents the 20 main use cases of the
Medication Safety Assistant system. The diagrams illustrate
the interactions between actors and the system's functional
capabilities. Detailed descriptions, basic flows, alternative
flows, and conditions are provided in the
Use-Case Specification Document.


% Reusable use-case figure
\newcommand{\usecasefigure}[3]{%
    \begin{figure}[!htbp]
        \centering
        \includegraphics[
            width=0.95\linewidth,
            height=0.37\textheight,
            keepaspectratio
        ]{images/UC-#1.png}
        \caption[UC-#1: #2]{%
            \textbf{UC-#1: #2.} #3%
        }
        \label{fig:uc-#1}
    \end{figure}
}


% UC-01
\usecasefigure{01}{Sign In}{
    Allows registered users to securely access the system
    according to their assigned roles.
}

% UC-02
\usecasefigure{02}{Enter and Confirm Medication Information}{
    Allows patients to enter, review, and confirm their
    medication information before further processing.
}

% UC-03
\usecasefigure{03}{Check Medication Safety}{
    Allows patients to check potential medication risks,
    including drug interactions and other safety concerns.
}

% UC-04
\usecasefigure{04}{Arrange HCP Review}{
    Supports authorized HCP review requests, case
    assignment, and review status tracking.
}

% UC-05
\usecasefigure{05}{Review Medication Safety Case}{
    Allows healthcare professionals to review authorized
    safety cases and communicate with patients.
}

% UC-06
\usecasefigure{06}{Set Up Medication Reminders}{
    Allows patients to set up and confirm medication reminders
    based on verified schedules, with HCP review when required.
}

% UC-07
\usecasefigure{07}{Manage Medication Reminders}{
    Allows patients to adjust notification settings and
    propose medication timing changes subject to HCP review.
}

% UC-08
\usecasefigure{08}{Record Dose Status}{
    Allows patients to record and confirm their
    medication-taking status.
}

% UC-09
\usecasefigure{09}{View Adherence History}{
    Allows patients to view, filter, and summarize
    their medication adherence history.
}

% UC-10
\usecasefigure{10}{View Medication and Safety History}{
    Allows patients to view medication records, examine past
    safety assessments, and compare assessment results.
}

% UC-11
\usecasefigure{11}{Ask a Medication Question}{
    Allows patients to ask medication-related questions,
    view chatbot conversation history, and access supporting references.
}

% UC-12
\usecasefigure{12}{Generate Patient Health Summary}{
    Allows patients to generate, review, and confirm
    a structured summary of their health information.
}

% UC-13
\usecasefigure{13}{Share Patient Health Summary with HCP}{
    Allows patients to share health summaries with authorized HCPs,
    confirm sharing permissions, and revoke access when needed.
}

% UC-14
\usecasefigure{14}{Review Medication Reference Data}{
    Allows dataset reviewers to verify reference data against
    supporting evidence, record quality decisions, and review history.
}

% UC-15
\usecasefigure{15}{Manage User Accounts and Roles}{
    Allows administrators to create and deactivate accounts,
    assign roles, reset passwords, and view account activity history.
}

% UC-16
\usecasefigure{16}{Report Technical Issue}{
    Allows registered users to report technical issues,
    review attachments, reply to support tickets,
    and track ticket status.
}

% UC-17
\usecasefigure{17}{Handle Technical Support Request}{
    Allows technical support staff to view assigned tickets,
    investigate issues, request additional information,
    and manage ticket resolution.
}

% UC-18
\usecasefigure{18}{Manage Patient Health Profile}{
    Allows patients to manage medical conditions, allergies,
    medication and adverse reaction histories,
    lifestyle information, and health measurements.
}

% UC-19
\usecasefigure{19}{Manage Personal Medication Cabinet}{
    Allows patients to add and update personal medicines,
    track usage status, view medication details and known
    side effects, and monitor expiration dates.
}

% UC-20
\usecasefigure{20}{Manage Personal Account}{
    Allows registered users to view and update their
    personal account information and change passwords.
}

\clearpage


% 4. Architecture diagrams; component responsibilities, services,
%    communications (HTTP/HTTPS/etc.); class diagrams for components.
%    Add \subsection{Component: <name>} for each component.
% ============================================================
% SECTION 4 -- LOGICAL VIEW (CONTENT-FIRST REVISION)
% ============================================================
% Insert as sections/04-Logical View.tex in the PA3 report.
% Requires: tabularx, longtable, array (already used by main.tex).
% No image file is required for this revision.
% Diagram insertion points are marked as comments for the next step.
% Class names and method signatures are PROPOSED logical interfaces,
% not a claim about implemented code or final database schema.

\section{Logical View}
\label{sec:logical-view}

The Medication Safety Assistant follows a \textbf{layered modular monolith
(three-tier architecture)}. The system separates the user interface,
application logic, and data access. This design keeps the five main business
modules independent in responsibility while allowing them to run inside one
Python backend application. The nine logical components described below are
\emph{responsibility groups}, not nine independent microservices.

\noindent\textbf{Architectural style and technologies.}\par\smallskip
\begingroup\small\renewcommand{\arraystretch}{1.15}
\begin{tabularx}{\textwidth}{|L{2.85cm}|Y|L{2.55cm}|}
\hline
\textbf{Tier} & \textbf{Main responsibilities} & \textbf{Languages / means} \\ \hline
Presentation tier & Browser-based user interface; input forms, chatbot view,
notifications, role-based screens. & TypeScript; HTML; CSS \\ \hline
Application tier & API boundary, authentication checks, five business modules,
patient profile, reference-data review and technical support services. & Python; internal calls and DTOs \\ \hline
Data tier & Shared repositories and the existing versioned Core Database V4.
Stores authorized workflow, safety and history records. & PostgreSQL; SQL \\ \hline
\end{tabularx}
\endgroup

The frontend technology is described as \emph{TypeScript-based} because that
language is present in the project repository. Its exact UI framework and the
Python API framework must be confirmed from source code before implementation
contracts are finalized. \textbf{HTTPS with JSON} is the proposed browser--API
interface (local development may use HTTP); backend modules communicate through
\textbf{in-process Python calls and typed data transfer objects (DTOs)}. The
shared repository layer accesses PostgreSQL through parameterized SQL.
No separate network service is required between the five business modules.

\noindent\textbf{Overall component architecture.}\par\smallskip
The presentation/API boundary routes requests to the appropriate business
component. Module 1 passes confirmed entries to Module 2; Module 2 supplies
identified ingredients to Module 3. Module 3 may create a saved safety case
for Module 4. Module 5 checks verified schedule data and the safety rules
before activating a new reminder. Supporting components provide patient
health information, reference data, user permissions, and persistence.

% INSERT THE OVERALL COMPONENT DIAGRAM HERE LATER.
% Suggested file: images/logical-overview.pdf
% \begin{figure}[H]
%   \centering
%   \includegraphics[width=0.97\linewidth,keepaspectratio]{images/logical-overview.pdf}
%   \caption{Logical components, three tiers, service interfaces and communication means.}
%   \label{fig:logical-components}
% \end{figure}

\noindent\textbf{Component responsibilities and provided services.}\par\smallskip
\begingroup\small\renewcommand{\arraystretch}{1.14}
\begin{longtable}{|L{0.285\textwidth}|L{0.49\textwidth}|L{0.125\textwidth}|}
\hline
\textbf{Component} & \textbf{Responsibilities and services} & \textbf{Related UC} \\ \hline
\endfirsthead
\hline
\textbf{Component} & \textbf{Responsibilities and services (continued)} & \textbf{Related UC} \\ \hline
\endhead
Presentation and API Access & Receive user actions, apply access checks, and display
permitted responses. Delegate clinical logic to business modules. & UC-01--18, UC-20 \\ \hline
Medication Intake and Chatbot (Module 1) & Accept manual, OCR, and chatbot
input; clarify missing fields; require explicit confirmation; support bounded
medication questions. & UC-02, UC-11 \\ \hline
Medication Identification (Module 2) & Resolve confirmed product names to
supported medication and active-ingredient references; retain ambiguous or
unresolved identity. & UC-02, UC-03 \\ \hline
Safety Assessment and Decision Support (Module 3) & Evaluate supported
interaction and patient-specific risks; save warnings, source limitations,
and reminder-gate decisions. & UC-03, UC-10 \\ \hline
HCP Review and Patient Health Summary (Module 4) & Manage consent-bound
simulated reviews and case conversations; generate and share patient-reviewed
health summaries. & UC-04, 05, 12, 13 \\ \hline
Reminders and Adherence (Module 5) & Maintain eligible reminders and
notification settings; record dose status and adherence history; keep
unverified schedule proposals pending. & UC-06--09 \\ \hline
Patient Health Profile & Store source-aware conditions, allergies, medication
use status, past ADRs, family/lifestyle history and measurements. & UC-18 \\ \hline
Reference Data and Dataset Review & Provide local medication, ingredient and
DDI lookup; record versioned, evidence-linked dataset review decisions. & UC-14 \\ \hline
Accounts, Technical Support and Persistence & Provide account/RBAC,
technical-ticket, audit and shared repository functions. These are separate
internal responsibilities, not one all-purpose class. & UC-01, 15--17, 20 \\ \hline
\end{longtable}
\endgroup

\noindent\textbf{Relationships and communication.}\par\smallskip
\begingroup\footnotesize\renewcommand{\arraystretch}{1.12}
\begin{longtable}{|L{0.175\textwidth}|L{0.185\textwidth}|L{0.22\textwidth}|L{0.325\textwidth}|}
\hline
\textbf{Source} & \textbf{Target} & \textbf{Means} & \textbf{Information / service} \\ \hline
\endfirsthead
\hline
\textbf{Source} & \textbf{Target} & \textbf{Means} & \textbf{Information / service} \\ \hline
\endhead
Web UI & API boundary & HTTPS / JSON (planned) & Authorized requests,
results, validation feedback and role-specific views. \\ \hline
API boundary & Business and support services & In-process Python service call
& Route authenticated commands and queries; enforce access checks. \\ \hline
Intake (M1) & Identification (M2) & Python / confirmed-intake DTO &
Confirmed source text, medication fields and uncertainty. \\ \hline
Identification (M2) & Reference Data & Python / repository lookup &
Medication, ingredient and mapping evidence. \\ \hline
Identification (M2) & Safety (M3) & Python / identified-medication DTO &
Identified ingredients, matching status and provenance. \\ \hline
Safety (M3) & Patient Profile & Python / profile-query DTO &
Conditions, allergies and current/discontinued medication status. \\ \hline
Safety (M3) & HCP Review (M4) & Saved case reference / DTO &
Safety warnings and data gaps for authorized simulated review. \\ \hline
Reminders (M5) & HCP Review (M4) & Python / review-request DTO &
Unverified schedules and proposed timing changes. \\ \hline
HCP Review (M4) & Reminders (M5) & Recorded review result / DTO &
Schedule review status; trigger eligibility recheck, not auto-approval. \\ \hline
Reminders (M5) & Safety (M3) & Python / eligibility-gate call &
Allow, hold or report missing safety evidence before activation. \\ \hline
HCP Review (M4) & Profile and History & Python / permitted query &
Generate summaries and read patient-confirmed information. \\ \hline
All authorized services & Shared repositories & Python repository methods &
Read/write approved domain records and audit history. \\ \hline
Shared repositories & PostgreSQL Core DB V4 & SQL / PostgreSQL connection &
Transactions and versioned data under the existing schema contract. \\ \hline
Intake adapters & OCR / optional local LLM & Local library / loopback HTTP &
Text extraction or draft NLP output, never clinical clearance. \\ \hline
Reminders (M5) & Browser notification capability & Browser API &
Show eligible notifications when permission is available. \\ \hline
\end{longtable}
\endgroup

\noindent\textbf{Architectural rules and boundaries.}\par\smallskip
Medication details must be \textbf{explicitly confirmed} before identification
and safety checks. Module 2 must not treat ambiguous product matches as verified.
Module 3 checks supported drug--drug interactions, duplicate ingredients,
drug--disease risks, and available contraindication/allergy evidence; missing
coverage is labeled \emph{Unassessed}, not ``safe''. Medication history
marked as discontinued is not automatically treated as currently used.
For severe interactions or other defined review conditions, \textbf{affected
new reminders remain on hold}. An HCP review is simulated and requires valid
patient sharing permission. It cannot prescribe, stop or automatically change
treatment; any eligible reminder activation still requires confirmed schedule
information and a safety recheck. Patient health entries preserve source,
status and uncertainty; patient confirmation is not clinical validation.
All persistence must respect the approved Core DB V4 contract. UC-19 (Personal
Medication Cabinet) is documented as a future feature and is excluded from
this architecture's current implementation scope.

% The following nine subsections follow the instructor's component template.
% Add the class diagram within each component subsection after class design.

\subsection{Component: Presentation and API Access}
\label{subsec:logical-presentation}
The browser UI provides screens for patients, HCPs, administrators, dataset
reviewers and technical support. The API boundary validates incoming requests,
checks authenticated identity and resource permissions, and invokes the correct
application service. Domain safety rules are not implemented in the UI.

\textbf{Provided services.}\par
Role-scoped screens, request validation, authorization,
service routing and result display. The UI is planned in TypeScript with HTML
and CSS; backend API endpoints are planned in Python.

\textbf{Key classes/interfaces for later design.}\par\smallskip
\begingroup\small\renewcommand{\arraystretch}{1.10}
\begin{tabularx}{\textwidth}{|L{0.32\textwidth}|Y|}
\hline \textbf{Candidate class} & \textbf{Responsibility} \\ \hline
\texttt{WebClient} & Submit UI commands; display role-permitted responses. \\ \hline
\texttt{ApiController} & Validate and route incoming API requests. \\ \hline
\texttt{AccessGuard} & Apply RBAC and resource ownership rules. \\ \hline
\texttt{SessionContext} & Provide the current identity, role and session scope. \\ \hline
\end{tabularx}\endgroup

% CLASS DIAGRAM TO BE INSERTED HERE AFTER DESIGN: logical-presentation
% Use a diagram with main attributes, operations and relationships.


\subsection{Component: Medication Intake and Chatbot (Module 1)}
\label{subsec:logical-intake}
This component implements the patient intake workflow. It accepts manual text,
OCR output and chatbot messages, separates draft candidates from confirmed
records, asks for missing information, and preserves negation, uncertainty and
source text. It routes medication questions without treating every drug mention
as a request to save medication. Optional OCR and local language-model adapters
provide extraction assistance only.

\textbf{Provided services.}\par
Start intake, extract candidate fields, clarify,
confirm or cancel a draft, save confirmed intake, and answer supported
medication questions with clear limitations. Confirmed intake is handed to
Module 2 through a typed DTO.

\textbf{Key classes/interfaces for later design.}\par\smallskip
\begingroup\small\renewcommand{\arraystretch}{1.10}
\begin{tabularx}{\textwidth}{|L{0.32\textwidth}|Y|}
\hline \textbf{Candidate class} & \textbf{Responsibility} \\ \hline
\texttt{IntentRouter} & Classify intake, safety-request and Q\&A intents. \\ \hline
\texttt{DialogueManager} & Manage clarification and explicit confirmation state. \\ \hline
\texttt{IntakeValidator} & Check missing fields, spans, negation and uncertainty. \\ \hline
\texttt{IntakeRepository} & Persist draft/confirmed intake using the approved contract. \\ \hline
\texttt{IntakeHandoff} & Send confirmed intake DTOs to Module 2. \\ \hline
\end{tabularx}\endgroup

% CLASS DIAGRAM TO BE INSERTED HERE AFTER DESIGN: logical-intake
% Use a diagram with main attributes, operations and relationships.


\subsection{Component: Medication Identification (Module 2)}
\label{subsec:logical-identification}
Module 2 receives confirmed intake and matches medication names to supported
source-scoped products and active ingredients. It separates matching confidence
from user confirmation and records unresolved identities. Only sufficiently
identified entries are passed to the supported risk checks; ambiguous entries
remain visible for clarification or re-confirmation.

\textbf{Provided services.}\par
Medication matching, ingredient resolution,
identification status and evidence/provenance retrieval. It reads the approved
reference data through the shared catalog, not a live external service.

\textbf{Key classes/interfaces for later design.}\par\smallskip
\begingroup\small\renewcommand{\arraystretch}{1.10}
\begin{tabularx}{\textwidth}{|L{0.32\textwidth}|Y|}
\hline \textbf{Candidate class} & \textbf{Responsibility} \\ \hline
\texttt{IdentificationService} & Coordinate medication and ingredient identification. \\ \hline
\texttt{MedicationMatcher} & Retrieve candidate product matches with evidence. \\ \hline
\texttt{IngredientResolver} & Resolve supported active ingredients. \\ \hline
\texttt{IdentificationResult} & Carry identity status, uncertainty and source IDs. \\ \hline
\end{tabularx}\endgroup

% CLASS DIAGRAM TO BE INSERTED HERE AFTER DESIGN: logical-identification
% Use a diagram with main attributes, operations and relationships.


\subsection{Component: Safety Assessment and Decision Support (Module 3)}
\label{subsec:logical-safety}
Module 3 analyzes identified currently used medications using approved local
reference data and available patient profile information. It handles supported
drug--drug interactions, duplicate active ingredients, drug--disease risks,
and contraindication/allergy risks. The result includes warnings, evidence,
source versions and limitations. It must not claim clinical certainty where
supporting information is missing.

\textbf{Provided services.}\par
Run a safety assessment, record results and
warnings, show an \emph{Unassessed} status when needed, and evaluate eligibility
for \emph{new} reminders. A severe interaction holds the affected new reminder
pending the appropriate review and recheck; it does not automatically stop
existing reminders or alter treatment.

\textbf{Key classes/interfaces for later design.}\par\smallskip
\begingroup\small\renewcommand{\arraystretch}{1.10}
\begin{tabularx}{\textwidth}{|L{0.32\textwidth}|Y|}
\hline \textbf{Candidate class} & \textbf{Responsibility} \\ \hline
\texttt{SafetyAssessmentService} & Coordinate supported risk checks and evidence. \\ \hline
\texttt{InteractionChecker} & Evaluate DDI and duplicate-ingredient risks. \\ \hline
\texttt{PatientRiskChecker} & Check supported patient-specific risks. \\ \hline
\texttt{SafetyGate} & Evaluate whether an affected new reminder can activate. \\ \hline
\texttt{SafetyReport} & Record warnings, status, limitations and sources. \\ \hline
\end{tabularx}\endgroup

% CLASS DIAGRAM TO BE INSERTED HERE AFTER DESIGN: logical-safety
% Use a diagram with main attributes, operations and relationships.


\subsection{Component: HCP Review and Patient Health Summary (Module 4)}
\label{subsec:logical-review}
Module 4 handles two authorized review purposes: a saved safety case and a
pending schedule verification/change request. It records sharing permission,
case assignment, conversations and simulated HCP outcomes. Separately, it
prepares a patient-reviewed health summary and manages summary-sharing access.
The two sharing permissions are independent: sharing a summary does not grant
access to every safety case.

\textbf{Provided services.}\par
Create/reuse a review request, authorize case
sharing, assign an eligible simulated HCP, record and view outcomes, prepare a
health summary, share/revoke a confirmed summary. A completed review can
request a fresh reminder eligibility check, but never directly modifies a
prescription.

\textbf{Key classes/interfaces for later design.}\par\smallskip
\begingroup\small\renewcommand{\arraystretch}{1.10}
\begin{tabularx}{\textwidth}{|L{0.32\textwidth}|Y|}
\hline \textbf{Candidate class} & \textbf{Responsibility} \\ \hline
\texttt{ReviewCaseService} & Create, assign and track review requests. \\ \hline
\texttt{SharingPermissionService} & Manage consent for cases and separately for summaries. \\ \hline
\texttt{ReviewConversation} & Store authorized HCP/patient case communication. \\ \hline
\texttt{ReviewOutcome} & Capture a traceable simulated HCP outcome. \\ \hline
\texttt{HealthSummaryService} & Prepare, confirm and share permitted health summaries. \\ \hline
\end{tabularx}\endgroup

% CLASS DIAGRAM TO BE INSERTED HERE AFTER DESIGN: logical-review
% Use a diagram with main attributes, operations and relationships.


\subsection{Component: Reminders and Adherence (Module 5)}
\label{subsec:logical-reminders}
Module 5 uses only a verified existing prescription or recorded HCP schedule to
create eligible medication reminders. It checks the safety gate before a new
reminder is activated. Patients may change notification settings without
changing prescribed timing. A proposed schedule change is stored as pending
review and does not replace the verified schedule unless the required
confirmation and safety checks are completed. The component also records dose
status and provides adherence history.

\textbf{Provided services.}\par
Set up an eligible reminder, manage notifications,
record a proposed schedule change, record Taken/Missed/Skipped statuses, and
show adherence records. Browser notifications are optional: an eligible
schedule remains viewable in the application when notifications are unavailable.

\textbf{Key classes/interfaces for later design.}\par\smallskip
\begingroup\small\renewcommand{\arraystretch}{1.10}
\begin{tabularx}{\textwidth}{|L{0.32\textwidth}|Y|}
\hline \textbf{Candidate class} & \textbf{Responsibility} \\ \hline
\texttt{ReminderService} & Create eligible reminders and manage notifications. \\ \hline
\texttt{ReminderEligibilityGate} & Check schedule verification and safety status. \\ \hline
\texttt{ScheduleChangeRequest} & Keep proposed timing changes pending review. \\ \hline
\texttt{DoseStatusService} & Save confirmed Taken/Missed/Skipped dose events. \\ \hline
\texttt{AdherenceHistoryService} & Return permitted historical adherence records. \\ \hline
\end{tabularx}\endgroup

% CLASS DIAGRAM TO BE INSERTED HERE AFTER DESIGN: logical-reminders
% Use a diagram with main attributes, operations and relationships.


\subsection{Component: Patient Health Profile}
\label{subsec:logical-profile}
This supporting component manages patient-entered medical conditions,
allergies, medication history (including current/discontinued status), previous
adverse drug reactions, family/lifestyle history and health measurements. The
records preserve available source, status, time and uncertainty. A patient's
confirmation validates the entry process, not the medical accuracy.

\textbf{Provided services.}\par
View, add and update permitted profile records;
return only relevant, source-aware patient context to safety assessment and
health-summary services. Data storage should reuse the existing Core DB V4
contract; unsupported categories may need logical adapters rather than new
database tables without a separate schema review.

\textbf{Key classes/interfaces for later design.}\par\smallskip
\begingroup\small\renewcommand{\arraystretch}{1.10}
\begin{tabularx}{\textwidth}{|L{0.32\textwidth}|Y|}
\hline \textbf{Candidate class} & \textbf{Responsibility} \\ \hline
\texttt{PatientProfileService} & View and update permitted patient information. \\ \hline
\texttt{HealthProfileEntry} & Carry sourced conditions, allergies or other health data. \\ \hline
\texttt{MedicationUseEntry} & Record current/discontinued use and provenance. \\ \hline
\texttt{ProfileValidator} & Validate entry completeness and uncertainty. \\ \hline
\end{tabularx}\endgroup

% CLASS DIAGRAM TO BE INSERTED HERE AFTER DESIGN: logical-profile
% Use a diagram with main attributes, operations and relationships.


\subsection{Component: Reference Data and Dataset Review}
\label{subsec:logical-reference}
This component exposes the approved local medication, ingredient and DDInter
mapping records to Module 2 and Module 3. It also supports the Dataset Reviewer
in comparing records with source evidence and saving data-quality decisions.
Review status and source provenance are distinct from independent clinical
validation.

\textbf{Provided services.}\par
Query versioned references, retrieve mapped DDI
evidence, review source records, and record Accept/Reject/Needs Correction
with traceable history. Runtime safety checks do not require calls to a live
external drug-data service.

\textbf{Key classes/interfaces for later design.}\par\smallskip
\begingroup\small\renewcommand{\arraystretch}{1.10}
\begin{tabularx}{\textwidth}{|L{0.32\textwidth}|Y|}
\hline \textbf{Candidate class} & \textbf{Responsibility} \\ \hline
\texttt{ReferenceCatalog} & Query versioned medication and ingredient records. \\ \hline
\texttt{DDIReferenceReader} & Read locally mapped interaction evidence. \\ \hline
\texttt{DatasetReviewService} & Review evidence and record data-quality decisions. \\ \hline
\texttt{ReviewDecision} & Represent a source-linked decision and reviewer status. \\ \hline
\end{tabularx}\endgroup

% CLASS DIAGRAM TO BE INSERTED HERE AFTER DESIGN: logical-reference
% Use a diagram with main attributes, operations and relationships.


\subsection{Component: Accounts, Technical Support and Persistence}
\label{subsec:logical-platform}
This component is a \emph{logical grouping of three internal services}, not a
single large class: (1) accounts and role-based access control, (2) technical
support tickets and status tracking, and (3) shared persistence and auditing.
Account operations cover sign-in, personal profile edits, authorized role
assignment and deactivation. Technical Support receives permitted diagnostic
information only, not unrestricted patient health records. The repository
layer provides bounded access to the existing Core DB V4.

\textbf{Provided services.}\par
Authenticate users, check roles and resource
access, manage allowed account actions, create/update support tickets, persist
permitted records and write audit history. Other modules call repository
interfaces rather than defining their own database schemas.

\textbf{Key classes/interfaces for later design.}\par\smallskip
\begingroup\small\renewcommand{\arraystretch}{1.10}
\begin{tabularx}{\textwidth}{|L{0.32\textwidth}|Y|}
\hline \textbf{Candidate class} & \textbf{Responsibility} \\ \hline
\texttt{AccountService} & Provide authentication and authorized account actions. \\ \hline
\texttt{RBACPolicy} & Check roles and resource access. \\ \hline
\texttt{SupportTicketService} & Create, assign and update technical issue tickets. \\ \hline
\texttt{CoreRepository} & Read and write through the existing DB V4 contract. \\ \hline
\texttt{AuditRecorder} & Keep actor, event time and change history. \\ \hline
\end{tabularx}\endgroup

% CLASS DIAGRAM TO BE INSERTED HERE AFTER DESIGN: logical-platform
% Use a diagram with main attributes, operations and relationships.


\noindent\textit{Design status.} This Logical View defines the component boundaries
and the candidate class responsibilities for subsequent UML class diagrams.
The methods, attributes and exact framework-specific classes will be refined
against the real implementation and Core DB V4. Class diagrams must show the
main attributes, operations and relationships required by the PA3 assignment.

% 5-6. The lecturer allows these sections to remain empty for PA3.
% Their headings are always shown so the document matches the template.
\section{Deployment}
\label{sec:deployment}
% Intentionally blank for PA3, as instructed in the RUP Word template.

\section{Implementation View}
\label{sec:implementation-view}
% Intentionally blank for PA3, as instructed in the RUP Word template.


\end{document}